\documentclass[11pt]{article}

\usepackage[a4paper,margin=1in]{geometry}
\usepackage{amsmath,amssymb,amsthm,mathtools}
\usepackage{mathrsfs}
\usepackage{enumitem}
\usepackage{hyperref}
\usepackage{tikz-cd}
\usepackage{array}

\hypersetup{
	colorlinks=true,
	linkcolor=blue,
	urlcolor=blue,
	citecolor=blue
}

\setlength{\parskip}{0.5em}
\setlength{\parindent}{0pt}

\newtheorem{theorem}{Theorem}[section]
\newtheorem{proposition}[theorem]{Proposition}
\newtheorem{lemma}[theorem]{Lemma}
\newtheorem{corollary}[theorem]{Corollary}
\theoremstyle{definition}
\newtheorem{definition}[theorem]{Definition}
\newtheorem{example}[theorem]{Example}
\newtheorem{exercise}[theorem]{Exercise}
\theoremstyle{remark}
\newtheorem{remark}[theorem]{Remark}

\newcommand{\R}{\mathbb{R}}
\newcommand{\C}{\mathbb{C}}
\newcommand{\Z}{\mathbb{Z}}
\newcommand{\Q}{\mathbb{Q}}
\newcommand{\F}{\mathbb{F}}
\newcommand{\A}{\mathbb{A}}
\newcommand{\PP}{\mathbb{P}}
\newcommand{\cO}{\mathcal{O}}
\newcommand{\cI}{\mathcal{I}}
\newcommand{\cL}{\mathcal{L}}
\newcommand{\Syz}{\mathrm{Syz}}
\newcommand{\im}{\mathrm{im}}
\newcommand{\kerR}{\mathrm{ker}}
\newcommand{\ord}{\mathrm{ord}}
\newcommand{\divv}{\mathrm{div}}
\newcommand{\Spec}{\mathrm{Spec}}
\newcommand{\Hom}{\mathrm{Hom}}
\newcommand{\rank}{\mathrm{rank}}
\newcommand{\Pic}{\mathrm{Pic}}
\newcommand{\Supp}{\mathrm{Supp}}
\newcommand{\Proj}{\mathrm{Proj}}
\newcommand{\Sym}{\mathrm{Sym}}

\title{Lecture Notes on Syzygies for Algebraic Curves\\
	\large From Sections to Relations to Geometry}
\author{}
\date{}

\begin{document}
	
	\maketitle
	
	\tableofcontents
	
	\section{Purpose of these notes}
	
	These notes explain \textbf{syzygies} in a way that is adapted to the study of algebraic curves and compact Riemann surfaces.
	
	The main goal is to understand the following progression:
	\[
	\text{divisor} \longrightarrow \text{line bundle} \longrightarrow \text{global sections}
	\longrightarrow \text{map to projective space} \longrightarrow \text{equations}
	\longrightarrow \text{syzygies}.
	\]
	
	So syzygies are not introduced in isolation. They arise naturally once a curve is embedded in projective space by sections of a line bundle.
	
	\section{The main slogan}
	
	The central slogan is:
	
	\begin{center}
		\emph{Generators tell you what you build from; syzygies tell you how the generators are tied together.}
	\end{center}
	
	For curves, the generators will often be global sections of a line bundle, or homogeneous equations defining the embedded curve. A syzygy is a relation among those generators.
	
	\section{First definition}
	
	\begin{definition}
		Let \(R\) be a ring and let \(f_1,\dots,f_r\in R\). A \textbf{syzygy} among \(f_1,\dots,f_r\) is a tuple
		\[
		(a_1,\dots,a_r)\in R^r
		\]
		such that
		\[
		a_1f_1+\cdots+a_rf_r=0.
		\]
	\end{definition}
	
	So a syzygy is a relation among the \(f_i\), with coefficients from the ring \(R\).
	
	\begin{remark}
		This looks like linear dependence, but it is richer than linear algebra because the coefficients \(a_i\) are not scalars from a field; they are elements of a ring.
	\end{remark}
	
	\section{The most basic example}
	
	Let
	\[
	R=k[x,y].
	\]
	Consider the generators
	\[
	f_1=x,\qquad f_2=y.
	\]
	Then
	\[
	(-y)x + x y = 0.
	\]
	So
	\[
	(-y,x)
	\]
	is a syzygy of \(x\) and \(y\).
	
	This is the simplest model to remember:
	\[
	\Syz(x,y)=R\cdot(-y,x).
	\]
	
	It already captures the meaning of syzygy: the generators \(x\) and \(y\) are not isolated; they are yoked together by a relation.
	
	\section{Why this matters for curves}
	
	Suppose \(X\) is a smooth projective curve over a field \(k\), or equivalently for intuition, a compact Riemann surface.
	
	To study \(X\), one often does the following:
	
	\begin{enumerate}[label=\arabic*.]
		\item choose a divisor \(D\) on \(X\),
		\item form the line bundle \(\cL=\cO_X(D)\),
		\item study the vector space \(H^0(X,\cL)\) of global sections,
		\item use a basis of \(H^0(X,\cL)\) to define a map
		\[
		X\to \PP^r,
		\]
		\item study the image by understanding its defining equations,
		\item then study the relations among those equations.
	\end{enumerate}
	
	Those last relations are syzygies.
	
	So syzygies arise naturally once one passes from intrinsic geometry on \(X\) to the projective geometry of its image.
	
	\section{From divisors to sections}
	
	Let \(X\) be a smooth projective curve. If \(D\) is a divisor on \(X\), then
	\[
	\cL=\cO_X(D)
	\]
	is the associated line bundle.
	
	The space
	\[
	H^0(X,\cO_X(D))
	\]
	consists of meromorphic functions \(f\) on \(X\) such that
	\[
	\divv(f)+D\ge 0,
	\]
	together with \(0\).
	
	Thus \(H^0(X,\cO_X(D))\) is a space of functions allowed to have poles bounded by \(D\).
	
	\begin{remark}
		This is the first place where your divisor intuition enters. The divisor controls which sections are allowed. Once we choose a basis of allowed sections, those sections may satisfy algebraic relations. Those relations eventually become equations and syzygies.
	\end{remark}
	
	\section{From sections to a projective map}
	
	Suppose
	\[
	V=H^0(X,\cL)
	\]
	has dimension \(r+1\), and choose a basis
	\[
	s_0,\dots,s_r.
	\]
	If these sections have no common zero, they define a morphism
	\[
	\phi_{\cL}:X\to \PP^r,\qquad
	p\mapsto [s_0(p):\cdots:s_r(p)].
	\]
	
	Now the image \(\phi_{\cL}(X)\subseteq \PP^r\) is a projective curve. We can ask:
	
	\begin{center}
		What homogeneous polynomials vanish on this image?
	\end{center}
	
	These homogeneous polynomials form the homogeneous ideal
	\[
	I_X \subseteq k[x_0,\dots,x_r].
	\]
	
	The generators of \(I_X\) are the equations of the embedded curve.  
	The syzygies among those generators measure how the equations themselves fit together.
	
	\section{The geometric meaning of syzygy on a curve}
	
	Suppose the ideal of the embedded curve is generated by homogeneous polynomials
	\[
	f_1,\dots,f_m\in S:=k[x_0,\dots,x_r].
	\]
	Then a syzygy is a tuple
	\[
	(a_1,\dots,a_m)\in S^m
	\]
	such that
	\[
	a_1f_1+\cdots+a_mf_m=0.
	\]
	
	Interpretation:
	
	\begin{itemize}[leftmargin=2em]
		\item the \(f_i\) are the visible equations cutting out the image of the curve,
		\item the syzygies are the hidden compatibility laws among these equations.
	\end{itemize}
	
	So equations tell us \emph{where the curve is}; syzygies tell us \emph{how the equations are organized}.
	
	\section{A first curve example: \(\PP^1\)}
	
	Let \(X=\PP^1\). Consider the line bundle
	\[
	\cO_{\PP^1}(2).
	\]
	Its global sections have basis
	\[
	s^2,\ st,\ t^2,
	\]
	where \([s:t]\) are homogeneous coordinates on \(\PP^1\).
	
	These sections define the map
	\[
	\PP^1\to \PP^2,\qquad
	[s:t]\mapsto [s^2:st:t^2].
	\]
	
	Let the homogeneous coordinates on \(\PP^2\) be \([x:y:z]\). Then the image satisfies
	\[
	y^2-xz=0.
	\]
	
	So the image is a conic.
	
	\subsection{What about syzygies?}
	
	Here the ideal is generated by a single polynomial:
	\[
	I=(y^2-xz).
	\]
	Since there is only one generator, there is no nontrivial first syzygy among the generators of the ideal.
	
	\begin{remark}
		This is important: a curve can have nontrivial defining equation but still have trivial first syzygies if its ideal is principal.
	\end{remark}
	
	\section{A better example: the twisted cubic}
	
	To really see syzygies geometrically, we should look at a curve whose ideal has several generators.
	
	Consider the map
	\[
	\PP^1\to \PP^3,\qquad
	[s:t]\mapsto [s^3:s^2t:st^2:t^3].
	\]
	This is the \textbf{twisted cubic}.
	
	Let the homogeneous coordinates on \(\PP^3\) be
	\[
	[x_0:x_1:x_2:x_3].
	\]
	The image is cut out by the three quadrics
	\[
	f_1=x_1^2-x_0x_2,\qquad
	f_2=x_1x_2-x_0x_3,\qquad
	f_3=x_2^2-x_1x_3.
	\]
	
	So the ideal is
	\[
	I=(f_1,f_2,f_3).
	\]
	
	Now syzygies appear.
	
	\subsection{Two explicit syzygies}
	
	One checks directly that
	\[
	x_2f_1 - x_1f_2 + x_0f_3 = 0
	\]
	and
	\[
	x_3f_1 - x_2f_2 + x_1f_3 = 0.
	\]
	
	So two syzygies are
	\[
	(x_2,-x_1,x_0)
	\qquad\text{and}\qquad
	(x_3,-x_2,x_1).
	\]
	
	These are genuine geometric syzygies: relations among the quadratic equations defining the curve.
	
	\subsection{Why this is penetrating}
	
	The twisted cubic is not just the common zero set of three quadrics. Those quadrics themselves are interlocked. The syzygies reveal that the ideal has internal structure.
	
	So syzygies are already measuring something deeper than merely listing equations.
	
	\section{The homogeneous coordinate ring viewpoint}
	
	Let
	\[
	S=k[x_0,\dots,x_r]
	\]
	and let \(I_X\subseteq S\) be the homogeneous ideal of the embedded curve \(X\subseteq \PP^r\). Then the homogeneous coordinate ring is
	\[
	R_X=S/I_X.
	\]
	
	To study \(R_X\), one often resolves it by free \(S\)-modules:
	\[
	\cdots \to F_2 \to F_1 \to F_0 \to R_X \to 0.
	\]
	
	Interpretation:
	
	\begin{itemize}[leftmargin=2em]
		\item \(F_0\) describes generators,
		\item \(F_1\) describes first syzygies,
		\item \(F_2\) describes syzygies among syzygies,
		\item etc.
	\end{itemize}
	
	This is the language of free resolutions.
	
	\section{Syzygies as kernels}
	
	Suppose \(I=(f_1,\dots,f_m)\subseteq S\). Then the generators define a surjective map
	\[
	\varphi:S^m\to I,\qquad
	(a_1,\dots,a_m)\mapsto a_1f_1+\cdots+a_mf_m.
	\]
	Its kernel is the module of syzygies:
	\[
	\Syz(f_1,\dots,f_m)=\kerR(\varphi).
	\]
	
	So syzygy is just another name for ``kernel of the map determined by the generators.''
	
	This is the clean algebraic core.
	
	\section{How this relates to line bundles}
	
	Now let us connect the projective picture back to the intrinsic geometry of the curve.
	
	Suppose \(\cL\) is a line bundle on \(X\), and let
	\[
	V=H^0(X,\cL).
	\]
	Then the sections of \(\cL\) generate higher tensor powers through multiplication maps
	\[
	\Sym^n V \longrightarrow H^0(X,\cL^{\otimes n}).
	\]
	
	These maps need not be injective or surjective.
	
	\subsection{Kernel as relation}
	
	An element in the kernel of
	\[
	\Sym^n V \to H^0(X,\cL^{\otimes n})
	\]
	is a homogeneous relation among degree-\(n\) monomials in the chosen sections.
	
	For example, in the \(\PP^1\) conic embedding using \(s^2,st,t^2\), we have
	\[
	(st)^2 - (s^2)(t^2)=0.
	\]
	This becomes
	\[
	y^2-xz=0.
	\]
	
	Thus the equations of the embedded curve come from kernels of multiplication maps among sections.
	
	This is already a very natural way to think:
	\[
	\text{relations among products of sections} \Rightarrow \text{projective equations}.
	\]
	
	Syzygies then arise as relations among these relations.
	
	\section{Canonical curves}
	
	One of the most important line bundles on a smooth projective curve \(X\) is the canonical bundle
	\[
	K_X.
	\]
	Its global sections are holomorphic differentials:
	\[
	H^0(X,K_X).
	\]
	If \(X\) is nonhyperelliptic, these sections define the canonical map
	\[
	X\to \PP^{g-1},
	\]
	where \(g\) is the genus.
	
	The image is called the \textbf{canonical curve}.
	
	The ideal of the canonical curve, and especially its syzygies, contain very fine information about the geometry of \(X\). This is one reason syzygies are central in the geometry of curves.
	
	\begin{remark}
		At this level, syzygies stop being merely formal algebra. They detect special linear series, gonality, and subtle projective properties of the curve.
	\end{remark}
	
	\section{A useful intuitive picture}
	
	Here is a layered picture especially suited to curves:
	
	\subsection*{Layer 1: divisor}
	A divisor tells you where zeros and poles are allowed.
	
	\subsection*{Layer 2: line bundle}
	The line bundle \(\cO_X(D)\) packages that divisor information geometrically.
	
	\subsection*{Layer 3: sections}
	Global sections are meromorphic objects respecting those pole bounds.
	
	\subsection*{Layer 4: projective coordinates}
	A basis of sections gives coordinates in projective space.
	
	\subsection*{Layer 5: equations}
	The image satisfies polynomial equations because the chosen sections are not algebraically independent.
	
	\subsection*{Layer 6: syzygies}
	Those equations themselves satisfy higher relations.
	
	So syzygies belong to the same hierarchy as divisors and sections; they simply live at a more hidden, more global, more algebraic layer.
	
	\section{A tiny intrinsic example of a relation}
	
	Let \(X=\PP^1\), and let
	\[
	u=s^2,\qquad v=st,\qquad w=t^2
	\]
	be sections of \(\cO_{\PP^1}(2)\).
	
	Then
	\[
	v^2=uw.
	\]
	This is a relation among products of sections.
	
	When we interpret \(u,v,w\) as projective coordinates \(x,y,z\), this becomes
	\[
	y^2-xz=0.
	\]
	
	So the conic equation comes from a relation among sections.
	
	This is a very good bridge from the divisor/line-bundle point of view to the algebraic notion of equations and syzygies.
	
	\section{Where higher syzygies enter}
	
	Suppose the ideal \(I_X\) of a curve is generated by \(f_1,\dots,f_m\). Then the first layer is:
	\[
	S^m \to I_X.
	\]
	Its kernel gives first syzygies.
	
	But that kernel is itself an \(S\)-module, so one may ask for generators of it. Then those generators may themselves satisfy relations. This gives second syzygies, then third syzygies, and so on.
	
	Thus one obtains a free resolution:
	\[
	\cdots \to F_2\to F_1\to F_0\to I_X\to 0
	\]
	or equivalently a resolution of \(S/I_X\).
	
	\begin{center}
		\emph{Syzygies are the successive hidden layers of structure behind an embedded curve.}
	\end{center}
	
	\section{Why the word ``syzygy'' is apt}
	
	Historically, ``syzygy'' means something like ``yoked together'' or ``aligned.''
	
	That is a good name because in a relation
	\[
	a_1f_1+\cdots+a_mf_m=0,
	\]
	the generators \(f_1,\dots,f_m\) are not independent. They are constrained to fit together in a specific way.
	
	For curves, the equations cutting out the image are yoked together by further compatibility laws. Those laws are syzygies.
	
	\section{Minimal free resolution: what it tells you}
	
	A minimal free resolution of the homogeneous coordinate ring of a projective curve records:
	
	\begin{itemize}[leftmargin=2em]
		\item how many generators the ideal has,
		\item in which degrees they occur,
		\item how many first syzygies there are,
		\item in which degrees those syzygies occur,
		\item and so on.
	\end{itemize}
	
	This numerical package is one of the sharpest algebraic fingerprints of the embedded curve.
	
	\section{Philosophical summary}
	
	For a curve \(X\), there are two viewpoints:
	
	\subsection*{Intrinsic viewpoint}
	Divisors, line bundles, meromorphic functions, differentials, cohomology.
	
	\subsection*{Extrinsic viewpoint}
	Embedding into projective space, equations, ideals, free resolutions, syzygies.
	
	Syzygies belong to the extrinsic viewpoint, but they arise from the intrinsic one because the embedding itself comes from sections of line bundles.
	
	Thus syzygies form a bridge:
	\[
	\text{intrinsic geometry of }X
	\quad\longleftrightarrow\quad
	\text{algebra of the projective embedding}.
	\]
	
	\section{Exercises}
	
	\begin{exercise}
		Let \(X=\PP^1\), and consider the map
		\[
		[s:t]\mapsto [s^2:st:t^2].
		\]
		Show directly that the image lies on the conic
		\[
		y^2-xz=0.
		\]
	\end{exercise}
	
	\begin{exercise}
		Explain why the ideal \((y^2-xz)\subseteq k[x,y,z]\) has no nontrivial first syzygies.
	\end{exercise}
	
	\begin{exercise}
		For the twisted cubic
		\[
		[s:t]\mapsto [s^3:s^2t:st^2:t^3],
		\]
		verify directly that
		\[
		x_2f_1-x_1f_2+x_0f_3=0
		\]
		where
		\[
		f_1=x_1^2-x_0x_2,\qquad
		f_2=x_1x_2-x_0x_3,\qquad
		f_3=x_2^2-x_1x_3.
		\]
	\end{exercise}
	
	\begin{exercise}
		Write down a second explicit syzygy among the generators \(f_1,f_2,f_3\) of the twisted cubic ideal.
	\end{exercise}
	
	\begin{exercise}
		Let \(D=2[\infty]\) on \(\PP^1\). Identify a basis of \(H^0(\PP^1,\cO(D))\), and explain how the relation among degree-\(2\) monomials in these sections gives the conic equation.
	\end{exercise}
	
	\begin{exercise}
		Why is it natural that relations among sections of a line bundle should produce equations in projective space?
	\end{exercise}
	
	\section{Solutions sketch}
	
	\subsection*{Exercise 1}
	Substitute
	\[
	x=s^2,\qquad y=st,\qquad z=t^2.
	\]
	Then
	\[
	y^2-xz=(st)^2-(s^2)(t^2)=0.
	\]
	
	\subsection*{Exercise 2}
	A first syzygy among generators of an ideal requires at least two generators. Since the ideal is principal, generated by one polynomial, there is no nonzero relation of the form \(a(y^2-xz)=0\) in the polynomial ring unless \(a=0\).
	
	\subsection*{Exercise 3}
	Expand:
	\[
	x_2(x_1^2-x_0x_2)-x_1(x_1x_2-x_0x_3)+x_0(x_2^2-x_1x_3)=0.
	\]
	All terms cancel.
	
	\subsection*{Exercise 4}
	The second syzygy is
	\[
	x_3f_1-x_2f_2+x_1f_3=0.
	\]
	
	\subsection*{Exercise 5}
	On \(\PP^1\), one may take basis \(1,z,z^2\) in the affine coordinate \(z\), or equivalently \(s^2,st,t^2\) homogeneously. The basic relation is
	\[
	(st)^2-(s^2)(t^2)=0,
	\]
	which becomes
	\[
	y^2-xz=0.
	\]
	
	\subsection*{Exercise 6}
	A basis of sections gives coordinate functions in projective space. Algebraic dependence among products of these sections becomes polynomial equations satisfied by the image.
	
	\section{Final conclusion}
	
	A syzygy is a relation among generators. For curves, this becomes especially meaningful after embedding the curve into projective space via sections of a line bundle.
	
	Then:
	
	\begin{itemize}[leftmargin=2em]
		\item sections give coordinates,
		\item relations among products of sections give equations,
		\item relations among equations give syzygies.
	\end{itemize}
	
	So syzygies are the hidden algebraic structure behind the visible projective image of the curve.
	
	They are important because they connect:
	\[
	\text{divisors and line bundles}
	\quad\text{to}\quad
	\text{equations and projective geometry}.
	\]
	
	That is why syzygies are central in the modern study of algebraic curves.
	
\end{document}